%% ============================================================================
%% Fichier  : annexes/Ann_B_outils.tex
%% Titre    : Référentiel des Outils
%% Version  : 0.1 -- novembre 2026
%% ============================================================================

\chapter{Référentiel des Outils}
\label{ann:outils}

Le cours ne demande qu'un tableur (Excel ou LibreOffice Calc) et l'outil de
calcul d'empreinte du système. Les formules sont données avec les noms
français des fonctions et le point-virgule comme séparateur.

\section*{Intégrité des fichiers}

\begin{tableaucours}{L{3.6cm}L{8.6cm}}{%
    \tableauHeader{Besoin} &
    \tableauHeader{Commande}
}
Empreinte sous Windows & \texttt{certutil -hashfile fichier.csv SHA256} \\
Empreinte sous Linux & \texttt{sha256sum fichier.csv} \\
Vérifier un manifeste & \texttt{sha256sum -c MANIFESTE\_SHA256.txt} \\
\end{tableaucours}

\section*{Formules de tableur}

\begin{tableaucours}{L{3.8cm}L{5.8cm}L{2.6cm}}{%
    \tableauHeader{Test} &
    \tableauHeader{Formule} &
    \tableauHeader{Chapitre}
}
Doublons sur une clé & \texttt{=NB.SI(\$F\$2:\$F\$9612;F2)} &
\ref{chap:prologue} \\
Tirage aléatoire & \texttt{=ALEA()} puis tri, valeurs figées &
\ref{chap:preuve} \\
Existence dans une autre liste & \texttt{=SIERREUR(RECHERCHEV(A2;
PERSONNEL!\$A:\$D;4;FAUX);"ABSENT")} & \ref{chap:acces} \\
Conflit de profils & \texttt{=SI(ET(C2=1;D2=1);"CONFLIT";"")} &
\ref{chap:acces} \\
Suite d'échecs & \texttt{=SI(B2="ECHEC";C1+1;0)} puis
\texttt{=MAX(C:C)} & \ref{chap:exploitation} \\
Seuil sans visa & \texttt{=SI(ET(F2>50000000;G2<>"OUI");"SANS VISA";"")} &
\ref{chap:applications} \\
Âge & \texttt{=DATEDIF(C2;DATE(2025;12;31);"y")} &
\ref{chap:donnees-fraude} \\
Premier chiffre & \texttt{=GAUCHE(TEXTE(A2;"0");1)*1} &
\ref{chap:donnees-fraude} \\
Proportion de Benford & \texttt{=LOG10(1+1/d)} & \ref{chap:donnees-fraude}
\\
\end{tableaucours}

\section*{NDAP SUKULU}

Les labs se font dans le bureau d'analyse de NDAP SUKULU, ouvert dans le
navigateur. Il contient un tableur, un éditeur de texte et un terminal.
Les pièces sont en lecture seule dans \texttt{preuves}~; on travaille sur
la copie du dossier \texttt{travail}.
